\section{田、役与水路：家户在大变局中怎样维持}
\label{sec:synthesis-work-land-water}

王朝的财力从不以“财政”二字直接落到田间。它要经过播种、收获、储藏、运输、簿册、差役和地方裁决，才会变成军粮、俸给或都城市场中的物资。两晋南北朝的战乱与迁徙不断打断这些环节：田主离开，佃者改投，水利失修，仓储被军队占用，劳动力被征发，税籍又可能跟不上实际住处。因而“江南开发”或“北方凋敝”若只作为时代结论，容易掩盖不同河谷、城镇和家庭承受的差异。

家户并不是国家账簿中的小格子。它也是照料老人和儿童、安排婚姻、储备种子、借贷、分工和应付临时征调的单位。一个户能否留在土地上，取决于水源、亲属、地主、里正、寺院与市场能否在危机中提供某种支持；同一张户籍所记录的“丁”或“田”则把这些复杂关系压缩成朝廷可征收的类别。理解这种不重合，才不会把制度上的“有田”“有户”误读为生活上的稳定。

\subsection{战争之后的土地不自动成为荒地}

西晋末年北方战乱造成城毁、人口死亡和逃散，这是年序材料足以确认的事实；然而，不能从“战乱”一词直接推断每一块土地都荒芜，或每一处幸存者都离开。有人可能在坞壁、山谷或近城处继续耕作，有人被军队征粮，有人以依附关系换取保护，有人转入新的政权。土地的占有、耕作者与课税者因而可能分离。后世史书更容易记录一座都城的陷落，较少记录一块田在第二年由谁播种，这正是综合叙事必须承认的视野限制。

关中和河北的政权更替尤其说明，粮食并不是只在国库中被争夺。前赵、后赵、前燕、前秦和北魏都必须在战役之外维持城郭、军队与官员；控制一处平原或河谷的意义，在于能否让耕作、征收和转运继续。前秦扩张时能够动员多地资源，淝水之后的政治崩解则使原先连通的征发网络迅速断裂。并非每一处因失败立即停耕，而是各地要重新面对谁能保护收成、谁有权征取、谁能提供种子与牲畜。

江南的情形也不能写成北方移民一到便开垦繁荣。长江下游原有孙吴以来的农业、水利、城镇与地方社会，北来人口、官员和军队进入后会带来新的需求、技术、竞争与税役。某些地区因人力、市场和政权保护而扩展，另一些地区则可能因战乱、洪水、征发或土地关系恶化而受损。用“开发”概括长期变化是可以的，但它不应取消本地居民、迁入者与依附者在不同条件下付出的劳动和代价。

\subsection{水道不是天然的财富}

长江、淮河、汉水和运河支流让南方朝廷具有不同于北方的运输条件。船运可以把重粮送往建康、荆州或前线，比陆路长途搬运更有效；它也依赖船只、船户、码头、堤防、征调和安全。水面宽阔并不会自动成为国家能力。若地方不愿提供船只，若战争控制了渡口，若粮食尚未从乡村汇集到仓，朝廷便无法把地理优势变成军事行动。

东晋和南朝北伐反复暴露这一点。前线指挥者所见的是河对岸的城池，后方家户所见的可能是被征走的船、粮与劳力。军队沿水道向北推进时，供给线要穿过多级地方官和市场；一旦战事拖延，原本用于维持城市与家庭的储备便可能转为军粮。史书中的“粮尽退兵”是一个军事结果，背后却有无数未被列名的收集、搬运、等待和失去。

隋的再统一常被视为大运河之前的整合起点，也应注意它继承的是前代早已存在的江淮与长江网络。五八九年灭陈之后，北方政权需要接收江南的州郡、仓储与船运，才能让军事胜利变成持续的赋役与调度。接收并非只更换官印：水路知识、地方船户、堤防维护和市场信用都不能凭诏书复制。再统一之所以要求长期行政工作，正在于这些看似日常的基础设施没有随着陈朝覆亡而自动归顺。

\begin{turningpoint}[北伐的粮运：胜利的代价先由沿线承担]
一场北伐可为朝廷带来合法性与短期领土，却要把粮、船、夫役和地方治安集中到前线。若收复区不能供给自身，南方腹地就要继续承担运输；若军队后撤，沿线居民又要面对新的征发或报复。军事行动不能只按将领的进退衡量，也必须问谁把它变成可移动的粮食。
\end{turningpoint}

\subsection{赋役如何改变家庭时间}

租、调、徭役和兵役对家户的影响，不只是交出多少物品。一次征调可能改变播种和收获的时间，迫使家中其他成员承担原本由壮丁完成的劳作；一次长途运输可能使债务和依附关系加深；一次迁居则可能使旧有田土、邻里担保和婚姻联系同时失效。传世制度材料多以国家应得的物资和劳力表述这些安排，较少保存家户如何凑齐。正因如此，不能把条文中的固定比例轻易变成某年某地的实际负担。

北魏均田三长制试图把土地、户籍和赋役接合起来，正是为了降低临时征取的不可预期性。若受田与纳税能在同一地方簿册中登记，国家便较容易计算，家户也较可能知道自己面对何种责任。然而，稳定的制度不等于轻松的制度。土地质量、迁徙、隐匿、地方强者的介入和官吏的能力都会改变实际负担；三长等基层中介既可帮助家户向上陈报，也可能拥有新的控制资源。

南朝的侨置与土断也把家庭时间拉入行政。寄居者何时改入实际住地，谁可以保留旧籍，哪一项役使由谁承担，都会影响家庭安排劳动与关系的方式。国家希望通过登记得到连续收入，地方社会则要在旧有保护和新责任之间协商。史书赞扬或批评改革者时常用道德与政绩的语言；从家户看，改革更像是一次对住处、身份和未来劳作的重新排序。

\subsection{市场、工匠与城市的胃口}

城市并不只消费乡村剩余，也组织工匠、运输者、商人和服务者。洛阳、长安、平城、邺、建康和江陵的宫城、寺院、军营与市场，都需要木材、砖瓦、金属、布帛、纸墨和食物。大规模营造能带来工作和交易，也会集中征发与资源。迁都洛阳、修建寺院或重整都城，因而既是政治展示，也改变附近道路、劳力和供应的方向。

工匠在材料中的身影常常比官员模糊。石窟、城墙和墓葬证明许多技艺曾被投入，题记偶尔留下施主或工匠的名字；我们却很难据此知道他们受雇、被役使还是如何迁移。不要以这种缺失为由把他们写成无脸的“劳动力”。更谨慎的写法是承认工程所需的组织规模，同时说明具体招募、工期与报酬往往不可确证。这样的保留使工匠不再被错误地归入某个统一的国家意志。

河西与西域的城镇市场展示另一种连接。绿洲农业、商旅、军队与宗教网络在此交会，文书中的借贷、交易或役使片段说明经济关系并不只由中央税制构成。它们不能替整个走廊画出完整经济图，却提醒我们：边地城市的生计依赖本地水源与远方道路的双重条件。政权更替会改变税收、保护与通行，却未必使市场和技能完全消失。

\begin{causalchain}[制度得失：把税役常规化，同时把风险写进家户]
登记、均田和地方编组能让国家较少依赖临时搜括，也可能让家户预期其责任；代价是迁徙者、无地者和未被正确登记者更容易落在制度之外。城市工程与军粮调度能够维持都城和边防，却把劳役、运输风险和收成波动转移给沿线家庭。国家能力的增长因此不能只以国库收入衡量，也要看谁承担了它的日常成本。
\end{causalchain}

\subsection{慢变量不替代事件}

土地、水路、家户和市场的变化很慢，却不能被写成脱离年序的背景。洛阳陷落使北方家庭的地权与劳役关系中断；东晋立国使江南重新安排户籍和供给；北魏统一北方后，均田三长为田与役提出新的连接；六镇起事又把军户、流民和城镇推入新的分布；侯景之乱、江陵失陷与陈亡分别重排南方城市的存粮、船运和人力。慢变量就在这些事件中改变速度，而事件也因慢变量的压力获得后果。

读者不必假装能够从有限材料看见每一户的账本。更重要的是拒绝把家庭劳动从国家史中删去：没有耕作、运输、照料、交易和维修，军队、寺院、都城与官署都无法持续。两晋南北朝的重新接合，最终发生在这些反复被登记、征调、迁徙又重新安顿的日常活动中。
